\subsection{Physical motivation for log domain modeling}

Delay spread (DS) and angular spread (AS) are nonnegative, often
zero inflated, and heavy tailed. Their empirical distributions are commonly
modeled in the logarithmic domain in standard channel models
\cite{tr38901,winner2}. For a spread variable $y$, define:
\begin{equation}
z(x) = \log(1 + y(x)),\qquad \hat y(x) = \exp(\hat z(x)) - 1
\label{eq:log_transform_spread}
\end{equation}

For two pixels with native spread values \(y_a\) and \(y_b\), their difference
in transformed space is
\begin{equation}
z_a-z_b
=
\log\!\left(\frac{1+y_a}{1+y_b}\right).
\label{eq:log_ratio_spread}
\end{equation}
Thus, a fixed error in \(z\) corresponds approximately to a fixed
multiplicative error in \(1+y\). This is more appropriate than native space
least squares for spreads because large NLoS tails would otherwise dominate
the fit and the near zero LoS mass would be ignored. The calibrated spread
prior is still a deterministic log domain ridge estimator; it is not claiming
to estimate the full spike plus tail probability distribution.

\subsection{Pixel wise notation for the spread prior fit}

Let \(n\) index a valid ground pixel from a training map. For a selected
metric \(m \in \{\mathrm{DS},\mathrm{AS}\}\), define the native target
\(y_{m,n}\), the transformed target \(z_{m,n}=\log(1+y_{m,n})\), and the
raw log prior value \(z^{(0)}_{m,n}\). The calibrated prior is linear in a
23-dimensional feature vector:
\begin{equation}
\hat z_{m,n}
=
\textbf{x}_{m,n}^{\top}\textbf{w}_{m,r(n)},
\qquad
\textbf{x}_{m,n}\in\mathbb{R}^{23},
\qquad
\textbf{w}_{m,r}\in\mathbb{R}^{23}.
\label{eq:spread_pixel_linear}
\end{equation}
The regime \(r(n)\) is a discrete key containing the metric, topology class,
LoS/NLoS state and UAV height bin. In other words, the spread prior does not fit one
global line through all pixels. It fits a family of small linear models, each
one specialized to a propagation regime.

\subsection{Raw spread prior in log domain}

\begin{equation}
\begin{aligned}
z^{(0)}(x)
={}&
\log\!\left(1+b^{\mathrm{native}}_{\mathrm{LoS/NLoS}}(x)\right)
+ b_{\mathrm{topology}}
+ a_1\log(1+d_{2D})
+ a_2\theta_{\mathrm{inv}}
\\
&+
a_3\rho_{41}
+ a_4 h_{41}
+ a_5 n_{41}
+ a_6 n_{41}\theta_{\mathrm{inv}} .
\end{aligned}
\label{eq:raw_spread}
\end{equation}

where $\theta_{\mathrm{inv}}(x) = 1 - \theta_{\mathrm{norm}}(x)$ (high value at low elevation). Example coefficients for delay spread: the native anchors are $b^{\mathrm{native}}_{\mathrm{LoS}} = 4.0$ ns and $b^{\mathrm{native}}_{\mathrm{NLoS}} = 11.0$ ns, which the implementation converts to $\log(1+4.0)$ and $\log(1+11.0)$ before adding the remaining log domain terms. The delay coefficients are $a_1 = 0.11$, $a_2 = 0.62$, $a_3 = 0.48$, $a_4 = 0.32$, $a_5 = 0.86$, $a_6 = 0.60$. After \eqref{eq:raw_spread}, the raw log prior is clipped to $[0,8]$, inverse transformed with $\exp(\cdot)-1$, and clipped to the metric upper bound. For reproducibility, the \(z^{(0)}(x)\) entries used later in \(\textbf{x}_{\mathrm{SP}}\) denote the log value recovered from this clipped native raw prior, i.e.,
\(\log(1+\mathrm{raw\_prior}_{\mathrm{clipped}}(x))\), not an unconstrained
copy of \eqref{eq:raw_spread}.

The raw prior is not a 3GPP formula copied directly into CKM. Instead, it
takes the 3GPP/WINNER lesson that spread statistics should be separated by
propagation state and represented in log space \cite{tr38901,winner2}, then
adds CKM specific map morphology terms. UAV measurement studies also motivate
the elevation angle dependence: low elevation links usually interact with more
rooftops and street canyons than near overhead links \cite{cai2019,izydorczyk2019}.
For reproducibility, Table~\ref{tab:spread_raw_constants} lists the fixed
raw prior constants used before the regime wise ridge calibration. These
values are fixed hyperparameters of \(z^{(0)}\). They are not the fitted ridge
weights \(\textbf{w}_{m,r}\), which are learned from training city pixels and
stored in the calibration file.

\begin{table}[H]
\centering
\caption{Raw prior constants fixed before the ridge calibration in \eqref{eq:raw_spread}.}
\label{tab:spread_raw_constants}
\scriptsize
\setlength{\tabcolsep}{3pt}
\begin{tabularx}{\textwidth}{lccX}
\toprule
Quantity & Delay spread & Angular spread & Meaning \\
\midrule
\(b^{\mathrm{native}}_{\mathrm{LoS}}\) & \(4.0\) ns & \(3.0^\circ\) &
Native LoS anchor converted to \(\log(1+b)\). \\
\(b^{\mathrm{native}}_{\mathrm{NLoS}}\) & \(11.0\) ns & \(7.0^\circ\) &
Native NLoS anchor converted to \(\log(1+b)\). \\
\(a_1,\ldots,a_6\) &
\(0.11,0.62,0.48,0.32,0.86,0.60\) &
\(0.05,0.52,0.60,0.38,0.58,0.35\) &
Fixed weights multiplying \(\log(1+d_{2D})\), \(\theta_{\mathrm{inv}}\),
\(\rho_{41}\), \(h_{41}\), \(n_{41}\), and
\(n_{41}\theta_{\mathrm{inv}}\), respectively. \\
\bottomrule
\end{tabularx}
\end{table}

The topology offset \(b_{\mathrm{topology}}\) is also fixed before ridge
calibration. In the order open sparse low rise, open sparse vertical, mixed
compact low rise, mixed compact mid rise, dense block mid rise, and dense block
high rise, the delay spread offsets are
\(0.00, 0.18, 0.06, 0.22, 0.24, 0.34\). For angular spread, the corresponding
offsets are \(0.00, 0.09, 0.05, 0.16, 0.18, 0.24\). These are log domain
offsets, not native ns/degree offsets.

Every term in \eqref{eq:raw_spread} has a physical interpretation:
\begin{itemize}
\item \(b_{\mathrm{LoS/NLoS}}\) is the base spread level for the propagation
      state. NLoS starts higher because blocked paths arrive after additional
      reflections and diffractions.
\item \(b_{\mathrm{topology}}\) is a morphology offset. Dense and high rise
      topologies receive larger offsets because they contain more reflectors.
\item \(\log(1+d_{2D})\) gives a smooth range dependence. Longer links have
      more opportunity to encounter delayed or angularly separated paths.
\item \(\theta_{\mathrm{inv}}\) increases at low elevation angle, where the
      UAV to ground ray intersects more urban structure before reaching the
      receiver.
\item \(\rho_{41}\), \(h_{41}\), and \(n_{41}\) describe local building
      density, mean building height, and NLoS support in a \(41\times41\)
      window. These are direct map derived proxies for multipath richness.
\item \(n_{41}\theta_{\mathrm{inv}}\) is an interaction term. It activates
      when low elevation geometry and local obstruction occur together.
\end{itemize}

\subsection{Calibrated spread prior}

\subsubsection{Stage 2: 23-feature vector and design matrix}

\begingroup
\small
\begin{equation}
\textbf{x}_{\mathrm{SP}}(x) = \begin{bmatrix}
z^{(0)}(x)^2,\ z^{(0)}(x),\ \log(1+d_{2D}),\ \theta_{\mathrm{norm}},\ \theta_{\mathrm{inv}}\\
h_{\mathrm{norm}},\ h_{\mathrm{norm}}^2,\ \rho_{15},\ \rho_{41},\ h_{15},\ h_{41}\\
n_{15},\ n_{41},\ n_{41}\log(1+d_{2D}),\ \rho_{41}\theta_{\mathrm{inv}},\ b_{41},\ 1\\
c_{41},\ t_{41},\ \theta_{\mathrm{norm}}\rho_{41},\ c_{41}\theta_{\mathrm{inv}},\ t_{41}\rho_{41},\ t_{41}n_{41}
\end{bmatrix}^\top
\label{eq:spread_calibration_vector}
\end{equation}
\endgroup

For reading, the vector in \eqref{eq:spread_calibration_vector} can be grouped into five conceptual blocks:
the raw prior block \(\textbf{x}_{\mathrm{prior}}\), the
range/elevation/height block \(\textbf{x}_{\mathrm{geom}}\), the local building
block \(\textbf{x}_{\mathrm{morph}}\), one bias term, and the interaction block
\(\textbf{x}_{\mathrm{int}}\). The displayed order follows the implementation,
so these blocks are not all contiguous in the column vector. Table~\ref{tab:spread_vector_blocks}
states the membership explicitly.

\begin{table}[H]
\centering
\scriptsize
\setlength{\tabcolsep}{3pt}
\renewcommand{\arraystretch}{1.05}
\begin{tabularx}{\textwidth}{p{2.75cm} X p{2.55cm}}
\toprule
\textbf{Block} & \textbf{Entries from \(\textbf{x}_{\mathrm{SP}}(x)\)} & \textbf{Purpose} \\
\midrule
\(\textbf{x}_{\mathrm{prior}}\) &
\(z^{(0)}(x)^2,\ z^{(0)}(x)\) &
Raw prior anchor and curvature. \\
\(\textbf{x}_{\mathrm{geom}}\) &
\(\log(1+d_{2D}),\ \theta_{\mathrm{norm}},\ \theta_{\mathrm{inv}},\ h_{\mathrm{norm}},\ h_{\mathrm{norm}}^2\) &
Range, elevation and UAV height. \\
\(\textbf{x}_{\mathrm{morph}}\) &
\(\rho_{15},\rho_{41},h_{15},h_{41},n_{15},n_{41},b_{41},c_{41},t_{41}\) &
Local urban morphology and blockage. \\
Bias &
\(1\) &
Regime offset. \\
\(\textbf{x}_{\mathrm{int}}\) &
\(n_{41}\log(1+d_{2D}),\ \rho_{41}\theta_{\mathrm{inv}},\ \theta_{\mathrm{norm}}\rho_{41},\ c_{41}\theta_{\mathrm{inv}},\ t_{41}\rho_{41},\ t_{41}n_{41}\) &
Joint geometry and morphology effects. \\
\bottomrule
\end{tabularx}
\caption{Block membership of the 23 entries in the spread calibration vector.}
\label{tab:spread_vector_blocks}
\end{table}

\begin{table}[!htbp]
\centering
\scriptsize
\setlength{\tabcolsep}{2.5pt}
\renewcommand{\arraystretch}{0.88}
\begin{tabularx}{\textwidth}{p{2.7cm} p{3.3cm} X}
\toprule
\textbf{Term} & \textbf{Mathematical role} & \textbf{Physical interpretation} \\
\midrule
\(z^{(0)}(x)^2\) & Quadratic raw prior basis. & Corrects curvature in the first spread estimate. \\
\(z^{(0)}(x)\) & Linear raw prior basis. & Anchors the fit to the log domain prior. \\
\(\log(1+d_{2D})\) & Smooth distance basis. & Captures range growth in delay and angle support. \\
\(\theta_{\mathrm{norm}}\) & Normalized elevation angle. & Separates overhead from grazing links. \\
\(\theta_{\mathrm{inv}}\) & Inverse elevation angle. & Highlights low elevation links with stronger obstruction sensitivity \cite{khawaja_survey,tr38901}. \\
\(h_{\mathrm{norm}},h_{\mathrm{norm}}^2\) & Linear and quadratic UAV height basis. & Learns nonlinear height effects. \\
\(\rho_{15},\rho_{41}\) & Fine/coarse local building density. & Counts nearby reflectors and blockers. \\
\(h_{15},h_{41}\) & Fine/coarse density weighted height. & Distinguishes low rise from high rise clutter. \\
\(n_{15},n_{41}\) & Fine/coarse local NLoS support. & Measures nearby geometric shielding. \\
\(b_{41}\) & Blocker depth term. & Measures building protrusion above receiver height. \\
\(1\) & Bias term. & Gives each regime an additive offset. \\
\(c_{41}\) & Transmitter clearance term. & Measures UAV clearance above local buildings. \\
\(t_{41}\) & Taller building area fraction. & Measures local area occupied by buildings above the UAV. \\
\bottomrule
\end{tabularx}
\caption{Term by term inventory of the direct terms in the spread calibration vector.}
\label{tab:spread_direct_term_inventory}
\end{table}

The first block lets the model bend the raw prior quadratically. The geometry
block handles range and height. The morphology block measures nearby
scatterers, while the interaction block expresses effects such as density
mattering more at low elevation. Some terms are correlated: for instance,
\(z^{(0)}(x)^2\) is determined by \(z^{(0)}(x)\), and
\(t_{41}(x)n_{41}(x)\) is tied to both taller building and NLoS support
features. Individual fitted coefficients should therefore not be read as
isolated physical constants; the terms were kept because the calibrated prior
improved with them.

\begin{table}[!htbp]
\centering
\scriptsize
\setlength{\tabcolsep}{2.5pt}
\renewcommand{\arraystretch}{0.88}
\begin{tabularx}{\textwidth}{p{3.0cm} p{3.2cm} X}
\toprule
\textbf{Term} & \textbf{Mathematical role} & \textbf{Physical interpretation} \\
\midrule
\(n_{41}\log(1+d_{2D})\) & NLoS distance interaction. & Lets range act differently in blocked streets. \\
\(\rho_{41}\theta_{\mathrm{inv}}\) & Density elevation interaction. & Activates dense, low elevation morphology. \\
\(\theta_{\mathrm{norm}}\rho_{41}\) & Overhead density interaction. & Allows dense areas to change under high elevation. \\
\(c_{41}\theta_{\mathrm{inv}}\) & Clearance elevation interaction. & Couples rooftop clearance with grazing geometry. \\
\(t_{41}\rho_{41}\) & Tall building density interaction. & Highlights dense blocks with buildings above the UAV. \\
\(t_{41}n_{41}\) & Tall building NLoS interaction. & Activates where tall buildings and shielding coincide. \\
\bottomrule
\end{tabularx}
\caption{Interaction terms in the spread calibration vector.}
\label{tab:spread_interaction_inventory}
\end{table}

\begin{figure}[H]
\centering
\begin{tikzpicture}[
  >=Stealth,
  box/.style={draw, thick, rounded corners=3pt, align=center,
              minimum height=0.85cm, font=\footnotesize},
  group/.style={box, minimum width=2.9cm, fill=blue!10},
  fitbox/.style={box, minimum width=3.2cm, fill=green!12},
  arr/.style={->, thick}
]
\node[group] (prior) at (0,0) {Raw prior\\terms};
\node[group] (geom) at (3.4,0) {Range, elevation\\and height};
\node[group] (morph) at (6.8,0) {Density, height\\and blockage};
\node[group] (inter) at (10.2,0) {Interaction\\terms};
\node[fitbox] (vec) at (5.1,-1.8) {Feature vector\\$\textbf{x}_{\mathrm{SP}}(x)\in\mathbb{R}^{23}$};
\node[fitbox] (ridge) at (5.1,-3.6) {Regime ridge fit\\$\hat z=\textbf{x}_{\mathrm{SP}}^{\top}\textbf{w}^{\mathrm{SP}}_{r}$};
\node[fitbox, fill=orange!15] (native) at (5.1,-5.4) {Native prior\\$\hat y=\exp(\hat z)-1$};
\draw[arr] (prior) -- (vec);
\draw[arr] (geom) -- (vec);
\draw[arr] (morph) -- (vec);
\draw[arr] (inter) -- (vec);
\draw[arr] (vec) -- (ridge);
\draw[arr] (ridge) -- (native);
\end{tikzpicture}
\caption[Feature grouping for the calibrated spread prior.]%
{Feature grouping used by the calibrated spread prior.
The calibrated model is linear in \(\textbf{x}_{\mathrm{SP}}\), but the features include
nonlinear transforms and interactions selected from propagation intuition.}
\label{fig:spread_feature_groups}
\end{figure}

The new features beyond the attenuation calibration vector are listed below. The
shared morphology notation follows the path loss prior construction in
Sec.~\ref{sec:nlos_branch}. In particular, \(b(\cdot)\) is the building mask,
\(\rho_k\) is the building density box mean, \(h_k\) is the density weighted
building height box mean, \(n_k\) is the ground NLoS support box mean, and
\(\theta_{\mathrm{norm}}\) is the normalized elevation angle. The same
reflected padding \(\mathrm{BoxMean}_k\) operator is used here, and
\(H_{\mathrm{norm}}=\SI{90}{m}\) is the building height normalizer. This keeps
the spread prior vector self contained while preserving the common feature
definitions used by the path loss prior.

\paragraph{Normalised transmitter height.}
\begin{equation}
h_{\mathrm{norm}} =
\min\!\left\{1,\,
\frac{\log\!\left(1+\max(h_{\mathrm{tx}}-h_{\mathrm{rx}},0)\right)}{\log(401)}
\right\}
\label{eq:hnorm}
\end{equation}
This compresses UAV height on a log scale and saturates it in \([0,1]\), using
\(\log(401)\) as the scale associated with a 400 m height difference. Samples
above that relative height are assigned \(h_{\mathrm{norm}}=1\), matching the
implementation. This cap is intentional: the transmitter height distribution has
only a small high altitude tail, so after splitting samples into
topology specific experts there is not enough support above 400 m to estimate a
separate high altitude response reliably. The vector also includes
\(h_{\mathrm{norm}}^2\), which is an
intentionally correlated basis expansion that lets the linear model learn a
curved height effect.

\paragraph{Blocker depth.}
\begin{equation}
b_{41}(x) = \frac{\mathrm{BoxMean}_{41}\!\left(\max(\mathrm{topology}(\cdot)-h_{\mathrm{rx}},0)\,b(\cdot)\right)(x)}{H_{\mathrm{norm}}}
\label{eq:b41}
\end{equation}
Measures how much buildings in the $41\times41$ window protrude above the receiver height.
Only building pixels contribute because of the multiplier \(b(\cdot)\). A tall
building directly next to the receiver therefore increases \(b_{41}\), while
flat open ground leaves it at zero.

\paragraph{Transmitter clearance.}
\begin{equation}
c_{41}(x) = \frac{\mathrm{BoxMean}_{41}\!\left(\max(h_{\mathrm{tx}}-\mathrm{topology}(\cdot),0)\,b(\cdot)\right)(x)}{H_{\mathrm{norm}}}
\label{eq:c41}
\end{equation}
This is a density weighted clearance feature: the clearance above each building
pixel is averaged by \(\mathrm{BoxMean}_{41}\) over the full \(41\times41\)
window and then normalized by \(H_{\mathrm{norm}}\). Ground pixels contribute
zero through the multiplier \(b(\cdot)\). A large \(c_{41}\) therefore means
both that nearby buildings are present and that the transmitter is well above
them. A small value can mean either little local building coverage, or buildings
that reach close to or above the UAV height.

\paragraph{Taller building fraction.}
\begin{equation}
t_{41}(x) = \mathrm{BoxMean}_{41}\!\left(\mathbf{1}[\mathrm{topology}(\cdot)>h_{\mathrm{tx}}]\cdot b(\cdot)\right)(x)
\label{eq:t41}
\end{equation}
This is the fraction of the full \(41\times41\) window area occupied by
building pixels taller than the UAV, not the fraction of buildings after
conditioning on there being a building pixel. It is zero if the UAV is above
all local buildings or if the local window contains no buildings, and grows
when the link is likely to be embedded in a high rise canyon. It is especially
useful for distinguishing dense but low rise zones from dense high rise zones.

\paragraph{Six class topology definition.}
The spread prior uses the six topology labels produced by the final prior
code. Their thresholds and class names are defined in
\eqref{eq:topology_class_sp_support}. These labels are map level routing keys
for the spread prior ridge models; they are not additional per pixel
morphology features.

\subsubsection{Stage 3: regime wise ridge regression}

The regime $r$ for the spread prior is keyed by: metric (DS or AS), topology class (6 classes), propagation state (LoS or NLoS), antenna height bin, giving up to $2 \times 6 \times 2 \times 3 = 72$ regimes.
The antenna height bins use the same tertile style thresholds as
\eqref{eq:ant_bin_support}: \texttt{low\_ant} for
\(h_{\mathrm{tx}}\le\SI{58.12}{m}\), \texttt{mid\_ant} for
\(\SI{58.12}{m}<h_{\mathrm{tx}}\le\SI{103.85}{m}\), and
\texttt{high\_ant} above \SI{103.85}{m}. These calibration bins are
different from the broader height reporting bins used later in the results
tables.

Valid pixels are subsampled at rate $p = 0.02$ (2\%). The ridge regression:
\begin{equation}
\hat{\textbf{w}}^{\mathrm{SP}}_{r} = \left(\textbf{X}_r^\top \textbf{X}_r + \lambda \textbf{D}^\top\textbf{D}\right)^{-1}\textbf{X}_r^\top \textbf{z}_r,\quad \lambda = 10^{-3}
\label{eq:ridge}
\end{equation}

Here \(\textbf{X}_r\in\mathbb{R}^{N_r\times23}\) stacks one row
\(\textbf{x}_{\mathrm{SP},n}^{\top}\) per sampled pixel in regime \(r\), and
\(\textbf{z}_r\in\mathbb{R}^{N_r}\) stacks the corresponding transformed
targets. The coefficient vector
\(\hat{\textbf{w}}^{\mathrm{SP}}_{r}\) minimizes the regularized least squares problem
\begin{equation}
\hat{\textbf{w}}^{\mathrm{SP}}_{r}
=
\arg\min_{\textbf{w}\in\mathbb{R}^{23}}
\left\|
\textbf{X}_r\textbf{w}-\textbf{z}_r
\right\|_2^2
+
\lambda
\left\|
\textbf{D}\textbf{w}
\right\|_2^2 ,
\label{eq:spread_ridge_objective}
\end{equation}
where \(\textbf{D}=\mathrm{diag}(d_1,\ldots,d_{23})\), with \(d_{23}=0\)
and \(d_i=1\) for \(i\ne23\). Thus, the final interaction term
\(t_{41}n_{41}\), which is the last entry of \eqref{eq:spread_calibration_vector}, is left
unpenalized by the ridge term in the frozen calibration used by the final
model. This is an implementation detail of the stored artifact, not a separate
physical assumption. The regularizer is small, but it prevents unstable
coefficients when a regime contains few valid pixels or when two morphology
features are strongly correlated.

Term by term, \(\textbf{X}_r^\top\textbf{X}_r\) is the feature correlation
matrix for regime \(r\), \(\lambda\textbf{D}^\top\textbf{D}\) adds a small
diagonal stabilizer to every coefficient except the final
\(t_{41}n_{41}\) coefficient, and
\(\textbf{X}_r^\top\textbf{z}_r\) is the feature target
correlation vector in log spread space. This is ridge regression: it is the
ordinary least squares fit plus a penalty that discourages very large
coefficients. The penalty is important here because many morphology features
are related; for example, \(\rho_{41}\), \(h_{41}\), \(b_{41}\) and \(t_{41}\)
all increase in dense high rise areas.

Sufficient statistics are accumulated incrementally:
\begin{align}
\textbf{A}_r &\mathrel{+}= \textbf{X}_{r,s}^\top \textbf{X}_{r,s} \quad \forall s\\
\textbf{v}_r &\mathrel{+}= \textbf{X}_{r,s}^\top \textbf{z}_{r,s} \quad \forall s
\end{align}
so the full design matrix is never materialised in memory.
This is important in practice because the calibration artifacts are already
large even after compression into fitted coefficients and profiles. The path loss
LoS calibration JSON reused by the final model is approximately 200\,000 lines, and the
spread calibration must be rerun separately for delay spread and angular
spread whenever the feature set changes.

\paragraph{Fallback hierarchy.} During fitting, each sampled pixel updates the
exact accumulator and the four broader fallback accumulators listed below. A
coefficient vector is written only for accumulators with at least
$4 \times 23 = 92$ sampled pixels. During inference, the implementation tries
the same keys in this order and uses the first one present:
\begin{enumerate}
\item exact regime: (metric, topology class, LoS/NLoS, antenna bin),
\item same topology class + LoS/NLoS + \emph{all} antenna heights,
\item same topology class + \emph{all} LoS states + all antenna heights,
\item global topology + LoS/NLoS + all antenna heights,
\item fully global fallback: (\texttt{metric}, \texttt{global},
      \texttt{all\_los}, \texttt{all\_ant}), where \texttt{all\_los} is the
      implementation key for all visibility states, not LoS only pixels.
\end{enumerate}

\subsubsection{Stage 4: inverse transform and clipping}

\begin{equation}
\hat y(x) = \exp\!\left(\hat z(x)\right) - 1,\qquad
\hat z(x) = \textbf{x}_{\mathrm{SP}}(x)^\top
\hat{\textbf{w}}^{\mathrm{SP}}_{r}
\label{eq:inv_log}
\end{equation}

Regime dependent clipping is applied after the inverse transform:
\begin{equation}
\hat y_{\mathrm{clip}}(x)
=
\min\!\left(
y_{\max,r},
\max\!\left(0,\exp(\hat z(x))-1\right)
\right),
\label{eq:spread_clip}
\end{equation}
with delay spread clipped to $[0, 400]$~ns and angular spread clipped to
$[0, 15^\circ]$ in LoS or $[0, 90^\circ]$ in NLoS. The asymmetric angular
limit is deliberate: LoS angular spread is dominated by the direct path and is
physically closer to a near zero spike plus small reflections, whereas NLoS
angular spread can occupy a much wider support.

\subsection{Delay and angular spread prior pipeline}

The topology class RMSE tables for the calibrated spread priors are reported in
Chapter~\ref{chap:results}. This subsection keeps the final pipeline diagram
and the method level data flow used by both spread modalities.

\begin{figure}[H]
\centering
\resizebox{0.98\textwidth}{!}{%
\begin{tikzpicture}[
    >=Latex,
    font=\footnotesize,
    block/.style={draw, thick, fill=blue!10, rectangle, rounded corners,
      text width=2.75cm, minimum height=1.18cm, align=center},
    input/.style={block, fill=gray!20},
    branch/.style={block, fill=orange!10},
    note/.style={draw, dashed, fill=yellow!10, rounded corners,
      text width=3.25cm, align=left, font=\scriptsize},
    arr/.style={->, thick},
    alabel/.style={font=\scriptsize, fill=white, inner sep=1.2pt}
]
\node[input] (inputs) at (0,1.9) {Geometry, masks and morphology\\
$d_{2D}, \theta, h_{\mathrm{tx}}, \rho, h, n$\\$b_{41}, c_{41}, t_{41}$};
\node[block] (raw_prior) at (4.0,1.9) {Stage 1\\Raw log prior\\$z^{(0)}(x)$};
\node[block] (features) at (8.0,1.9) {Stage 2\\Feature vector\\$\textbf{x}_{\mathrm{SP}}(x)\in\mathbb{R}^{23}$};
\node[note] (keys) at (4.0,-2.1) {Regime \(r\) keys:\\metric (DS/AS), topology, LoS/NLoS, antenna bin};
\node[block] (ridge) at (8.0,-2.1) {Stage 3\\Regime ridge\\$\hat z=\textbf{x}_{\mathrm{SP}}^{\top}\hat{\textbf{w}}^{\mathrm{SP}}_{r}$};
\node[branch] (invert) at (12.0,-2.1) {Stage 4\\Inverse transform\\$\hat y=\exp(\hat z)-1$};
\node[input] (clip) at (16.0,-2.1) {Calibrated spread prior\\$\hat y(x)$ after clipping};
\draw[arr] (inputs) -- node[alabel, above=5pt, pos=0.55] {maps} (raw_prior);
\draw[arr] (raw_prior) -- node[alabel, above=5pt, pos=0.55] {$\mathbb{R}^{H\times W}$} (features);
\draw[arr] (features.south) -- node[alabel, right=7pt, pos=0.45] {$\mathbb{R}^{23\times H\times W}$} (ridge.north);
\draw[arr] (ridge) -- node[alabel, above=5pt, pos=0.55] {$\mathbb{R}^{H\times W}$} (invert);
\draw[arr] (invert) -- node[alabel, above=5pt, pos=0.55] {$\mathbb{R}^{H\times W}$} (clip);
\draw[arr, dashed] (keys) -- (ridge);
\draw[arr, dashed] (inputs.south) .. controls +(0,-1.35) and +(-2.3,-1.35) .. (features.south west);
\end{tikzpicture}%
}
\caption[Pipeline for the delay and angular spread priors.]%
{Pipeline for the delay and angular spread priors. Both spread modalities share this identical linear architecture, predicting expected spread in the log domain via regime wise ridge regression over 23 spatial and geometric features.}
\label{fig:spread_prior_pipeline}
\end{figure}
